\chapter{Pore analysis}

\section{Specimen preparation for SEM}

To obtain a usable pore analysis, a good specimen preparation is important.
If the specimen is mechanically polished, there is a danger of smearing the meaterial in the pores or introducing scratches hindering a good segmentation.

{\color{red} [...] }

\ac{ArBIB} also allows to cut cross sections of \ac{CSH}-needles of hydrated \ac{C3S} (see Figure~\ref{fig:csh-arbib}).
This reveals a star-like footprint.

\begin{figure}[htb!]
    \centering
	
	\begin{subfigure}{0.47\textwidth}
		\includegraphics[width=\textwidth]{methods/artistic-csh-anim/C3S_Bruchflaeche_101.jpg}
		\caption{Fractured surface of hydrated alite.}
	\end{subfigure}
	\hfill
	\begin{subfigure}{0.47\textwidth}
		\includegraphics[width=\textwidth]{methods/artistic-csh-anim/C3S_7d_cryoBIB_007.jpg}
		\caption{\ac{ArBIB} cross section of hydrated alite.}
	\end{subfigure}

	\begin{subfigure}{0.47\textwidth}
		\includegraphics[width=\textwidth]{methods/artistic-csh-anim/Grundflaeche.pdf}
		\caption{Footprint of \ac{CSH}-needles from hydrated alite.}
	\end{subfigure}
	\hfill
	\begin{subfigure}{0.47\textwidth}
		\ifdefined\PRINTABLE
			\includegraphics[width=\textwidth]{methods/artistic-csh-anim/still.png}
		\else
			\animategraphics[palindrome,autoplay,width=\textwidth]{12}{methods/artistic-csh-anim/Test00}{01}{30}
		\fi
		\caption{Artistic 3D reconstruction of \ac{CSH}.}
	\end{subfigure}
	
    \caption{Artistic 3D reconstruction of \ac{CSH}-fibres covering a \ac{C3S}-particle based on a \ac{ArBIB} cross section of hydrated alite. \cite{rossler.2023}}
    \label{fig:csh-arbib}
\end{figure}


\section{Pore analysis using LT-DSC, MIP, CT and SEM}

{\color{red} TODO: discuss the problems with the analysis?

New specimens are planned to measure with the new regime. Compare to the old regime.}


\color{black}